This section introduces the structure of core \GMR and \SOE models
included in the \GTAPINGAMS distribution. Both models are static, and
track the production and distribution of goods in the global economy.
In the \GTAP 12 Data Base the world is divided into 163 regions
(typically representing individual countries) and 65 commodities (or
goods), but computational constraints tend to limit the number of
regions and goods which can be included in a single
model. Datasets are easily aggregated which can be useful for
both empirical applications focusing on a subset of regions and for
debugging steps at the model formulation stage.

In each region, final demand structure is composed of public and
private expenditure across goods. Decisions about the allocation of
resources are decentralized, and the representation of behavior by
consumers and firms in the model follows the canonical microeconomic
optimization framework: (i) consumers maximize welfare subject to a
budget constraint with fixed levels of investment and public
expenditure; (ii) producers combine intermediate inputs, and primary
factors (several categories of labor, land, resources and physical
capital) at least cost for given technology.\footnote{Details about
  \GTAP sectors and primary factors are provided in Appendix
  \ref{sectors} of \DOCU.  Readers of this paper
  who are interested in using \GTAPINGAMS as a research tool are
  encouraged to refer to Appendix \ref{installation} of
  \DOCU for installation instructions. The
  computer programs described have been tested using data files from
  the \GTAP 12 GDX package, contained in
  \texttt{GDX12a\_AY16994.zip}. A license for these data files must be
  purchased separately by the users. In addition, a licensed
  installation of \GAMS (version 24.7.3 or higher) is required. }

While \GMR and \SOE models are based on the same data and employ the
same representation of production technology and consumer preferences,
they of course differ in their representation of trade. The \GMR model
is based on a standard Armington representation of bilateral trade
flows. An implication is that production and consumption decisions in
a given country (or group of countries) will affect world prices, and
the magnitude of this effect will mainly depend on the elasticity
estimates (a measure of country-level market-power). By contrast, the
\SOE model considers the case in which production and consumption
decision in a region country (or group of country) do not affect world
prices. This is modeled by treating supply and demand by the ``rest of
the world'' (\textsc{row}, i.e. countries that are not considered
relevant for the \SOE analysis) as perfectly elastic.  Corner
solutions are avoided through the assumption that output destined for
the domestic and export markets are differentiated products
\citep[see][]{DeMeloRobinson1989}.\footnote{In the present
implementation, international closure between the region of interest
and the \textsc{row} is achieved by fixing the value of the current
account at its benchmark level, and permitting the real exchange rate
to clear the market. Other assumptions are of course possible and are
left for future research. } Therefore, from the perspective of a given
region, the impact of policy evaluated in \GMR and \SOE models will
differ through their impact on the terms of trade.

In the following, we start by describing the basic notation, and then
present the structure of the data together with benchmark accounting
identities. We then present a ``primal'' description of agents'
optimization problems (i.e. specified in terms of quantity variables),
which leads to the equilibrium conditions presented in the subsequent
section.

\subsection{Notation}

The notation used in the model is summarized in the Tables \ref{sets}
- \ref{prices}. Table \ref{sets} defines the various dimensions which
characterize an instance of the model, including the set of
sectors/commodities ($i,j$), the set of regions ($r,s)$, the set of
factors of production ($f$). Set $g$ combines the production sectors
$i$ and private and public consumption demand (indices \texttt{"C"}
and \texttt{"G"}) and investment demand (index \texttt{"I"}). It
allows for a much tighter formulation of the model as they can all be
conceived of ``goods'' produced in similar fashion. To simplify the
exposition of the model, however, we describe private consumption,
public consumption and investment demand as stand alone components.

\renewcommand{\arraystretch}{1}
\renewcommand{\tabcolsep}{1mm}
\begin{table}[t!]
\begin{center}
\caption{Definitions of set indices.} \label{sets}
\small
\begin{tabular}{lp{4in}}
\toprule
Set &Definition  \\
\addlinespace
\midrule
\addlinespace
$i,j$ & Sectors, an aggregation of the 65 sectors in the \GTAP 12
Data Base \\
$g$ & Production sectors $i$, plus private consumption \texttt{"C"},
public demand \texttt{"G"} and investment \texttt{"I"} \\
$r,s$ & Regions, an aggregation of the 163 regions in the
\GTAP 12 Data Base\\
$f$ & Factors of production (consisting of
\textit{mobile factors}, $f \in m_f$, skilled labor (i. officials,
mangers and legislators (ISCO-88 Major Groups 1-2), ii. technicians
and associated professionals, iii.
clerks, and iv. service and market sales workers), unskilled labor,
capital, and \textit{sector-specific}, $f \in s_f$, agricultural land
and other resources)\\ \addlinespace \bottomrule
\end{tabular}
\end{center}
\vspace{0.4in}
\end{table}

Table \ref{activities} defines the primal variables (activity levels)
which characterize an equilibrium. The model determines values of all
the variables except international capital flows, a parameter which
would be determined endogenously in an intertemporal model. Table
\ref{activities} also displays the concordance between the variables
and their \GAMS equivalents.

\renewcommand{\arraystretch}{1}
\renewcommand{\tabcolsep}{1mm}
\begin{table}[t!]
\begin{center}
\caption{Definitions of activity levels (quantity variables).} \label{activities}
\small
\begin{tabular}{lllr}
\toprule
Variable &Definition  & GAMS variable & Benchmark (\GTAP) value \\
\addlinespace
\midrule
\addlinespace
$Y_{ir}$&       Production & \texttt{Y(i,r)} &\texttt{vom(i,r)}\\
\\
$C_{r}$&        Discretionary consumption & \texttt{Y("c",r)} &\texttt{vom("c",r)}\\
\\
$G_{r}$&        Aggregate public & \texttt{Y("g",r)}&\texttt{vom("g",r)} \\
\\
$I_{r}$&        Aggregate investment & \texttt{Y("i",r)}&\texttt{vom("i",r)} \\
\\
$M_{ir}$&       Aggregate imports & \texttt{M(i,r)}&\texttt{vim(i,r)}\\
\\
$X_{ir}$&       Trade flows with rest of world regions &
\texttt{X(i,r)}&\texttt{vem(i,r)}\\
\\
$\textit{FT}_{fr}$&      Factor transformation &
\texttt{FT(f,r)}&\texttt{evom(f,r)}\\
\\
$\textit{YT}_j$  &       International transport services & \texttt{YT(j)}&\texttt{vtw(j)}\\
\addlinespace
\bottomrule
\end{tabular}
\end{center}
\vspace{0.4in}
\end{table}

Table \ref{prices} defines the relative price variables for goods and
factors in the model. As is the case in any Shoven-Whalley \CGE model,
the equilibrium conditions determine \textit{relative} rather than
\textit{nominal} prices.

\renewcommand{\arraystretch}{1}
\renewcommand{\tabcolsep}{1mm}
\begin{table}[t!]
\begin{center}
\caption{Definitions of price variables.} \label{prices}
\small
\begin{tabular}{lp{3in}r}
\toprule
Variable &Definition & GAMS variable\\
\addlinespace
\midrule
\addlinespace
$p^C_{r}$      &       Final demand price index for the Cobb-Douglas
demand system &\texttt{P("C",r)}\\
\\
$p^G_{r}$      &       Public provision price index  &\texttt{P("G",r)}\\
\\
$p^I_{r}$      &       Investment price index &\texttt{P("I",r)}\\
\\
$p^Y_{ir}$     &       Supply price, gross of indirect producer
taxes (when $\eta^{DX}_{i} < \infty$, the supply price is
replaced by and export price, $p^X_{ir}$ and domestic supply price,
$p^D_{ir}$) &\texttt{P(i,r)}\\
\\
$p^M_{ir}$     &       Import price, gross of export
taxes and tariffs.&\texttt{PM(i,r)}\\
\\
$p^E_{ir}$     &       Export price from endogenous or exogenous
regions.&\texttt{PE(i,r)}\\
\\
$p^T_j$         &       Marginal cost of transport services
&\texttt{PT(j)}\\
\\
$p^F_{f,r}$     &       Price of mobile primary factors ($m_f$
include labor, land and resources) &\texttt{PF(f,r)}\\
\\
$p^S_{s_f,ir}$     &       Price of sector-specific primary
factors&\texttt{PS(sf,i,j)}\\
\addlinespace
\bottomrule
\end{tabular}
\end{center}
\vspace{0.4in}
\end{table}

While the core \GTAPINGAMS model is base on Cobb-Douglas final demand,
the canonical models incorporate logic for two additional demand
systems: the \LES and the \CDE demand system. In particular, the \LES
representation distinguishes between subsistence demand and
discretionary demand, and we define the additional variables in Table
\ref{dsysprices}. Similarly, we define an Armington composite price
index for products entering \CDE final demand.  (Details of these
demand systems are provided in Appendices \ref{CDE} and \ref{LES} of
\DOCU.)

\begin{table}[t!]
\begin{center}
\caption{Additional variables for \LES and \CDE demand systems.} \label{dsysprices}
\small
\begin{tabular}{lp{3in}r}
\toprule
Variable &Definition & GAMS variable\\
\addlinespace
\midrule
\addlinespace
$SD_{r}$&       Subsistence demand & \texttt{Y("sd",r)} \\
\\
$DD_{r}$&       Discretionary demand & \texttt{Y("dd",r)} \\
\\
$p^{DD}_{r}$   &       Discretionary demand price index &\texttt{P("dd",r)}\\
\\
$p^{SD}_{r}$   &       Subsistence demand price index &\texttt{P("sd",r)}\\
\\
$p^A_{iCr}$     &      Armington composite price representing the
market price of goods entering \CDE demand.&\texttt{PA(i,r)}\\
\addlinespace
\bottomrule
\end{tabular}
\end{center}
\vspace{0.4in}
\end{table}

Finally, Table \ref{taxes} reports the definition of tax and subsidy
rates applied in the model, both in terms of the notation employed to
describe the model and that used in the \GAMS code. Note that revenues
from taxes and subsidy expenditures do not appear as explicit
variables in the \GTAP Data Base and are defined on the basis of
expenditures and tax rates. We come back to this below.

\renewcommand{\arraystretch}{1}
\renewcommand{\tabcolsep}{1mm}
\begin{table}[t!]
\begin{center}
\caption{Tax and subsidy rates (net basis unless noted).} \label{taxes}
\small
\begin{tabular}{lp{6cm}lr}
\toprule
Parameter & Definition &  & \GAMS Parameter \\
\addlinespace
\midrule
\addlinespace
$t^o_{ir}$ & Output taxes (gross basis)& & \texttt{rto(i,r)}    \\
\\
$t^{f}_{fir}$   &Factor taxes &  & \texttt{rtf(f,i,r)}   \\
\\
$t^{fd}_{ijr}$ & Intermediate input taxes &Domestic&\texttt{rtfd(i,j,r)}      \\
$t^{fi}_{ijr}$ &                          &Imported& \texttt{rtfi(i,j,r)}      \\
\\
$t^{pd}_{ir}$ & Consumption taxes  & Domestic& \texttt{rtfd(i,"C",r)}     \\
$t^{pi}_{ir}$ &                 &Imported& \texttt{rtfi(i,"C",r)} \\
\\
$t^{gd}_{ir}$ & Public demand taxes &Domestic & \texttt{rtfd(i,"G",r)}   \\
$t^{gi}_{ir}$&                 &Imported & \texttt{rtfi(i,"G",r)} \\
\\
$t^{gd}_{ir}$ & Investment demand taxes &Domestic & \texttt{rtfd(i,"I",r)}   \\
$t^{gi}_{ir}$&                 &Imported& \texttt{rtfi(i,"I",r)} \\
\\
$t^{xs}_{isr}$  &Export subsidies&  & \texttt{rtxs(i,s,r)}   \\
\\
$t^{ms}_{isr}$  &Import tariffs& & \texttt{rtms(i,s,r)}\\
\addlinespace
\bottomrule
\end{tabular}
\end{center}
\vspace{0.4in}
\end{table}

\subsection[Accounting identities]{Benchmark data structure and accounting identities}

The economic structure underlying the \GTAP dataset and model is
illustrated in Figure \ref{fig:region}. Symbols in this flow chart
correspond to variables in the economic model (see Table
\ref{activities}): $Y_{ir}$ is the production of good $i$ in region
$r$, $C_r$, $I_r$ and $G_r$ portray private consumption, investment
and public demand, respectively, $M_{ir}$ portrays the import of good
$i$ into region $r$, $\textit{RA}_r$ stand for representative
consumers, and $FT_{s_{f}r}$ is the activity through which the set of
sector-specific factors of production ($s_f$) are allocated to
individual sectors. Further, solid lines represent commodity and
factor market flows, while dotted lines indicate tax revenues and
transfers.

\begin{sidewaysfigure}
\begin{center}
\includegraphics{figures/gtapflowchart.tikz}
\caption{Regional economic structure.} \label{fig:region}
\end{center}
\end{sidewaysfigure}

Domestic and imported goods markets are represented by horizontal
lines at the top of the figure. Domestic production
($\textit{vom}_{ir}$) is distributed to exports
($\textit{vxmd}_{irs}$), international transportation services
($\textit{vst}_{ir}$), intermediate demand ($\textit{vdfm}_{ijr}$),
household consumption ($\textit{vdfm}_{iCr}$), investment
($\textit{vdfm}_{iIr}$), and government consumption
($\textit{vdfm}_{iGr}$).\footnote{Recall that in the \GAMS
implementation of the model the index \texttt{"g"} includes all
sectors represented in the model plus private consumption
\texttt{"C"}, public demand \texttt{"G"} and investment \texttt{"I"}.
(See Table \ref{sets}.) For the LES demand representation, it also
includes price indices for discretionary demand \texttt{"dd"} and
subsistence demand \texttt{"sd"}. } The accounting identity in the
\GTAP 12 dataset is thus:
$$\textit{vom}_{ir} = \sum_s \textit{vxmd}_{irs} + \textit{vst}_{ir}
+ \sum_j \textit{vdfm}_{ijr} + \textit{vdfm}_{iCr} +
\textit{vdfm}_{iIr} + \textit{vdfm}_{iGr}\,,$$ where $j$ indexes all
goods. Similarly, imported goods (with aggregate value
$\textit{vim}_{ir}$) enter intermediate demand
($\textit{vifm}_{ijr}$), private consumption ($\textit{vifm}_{iCr} +
\textit{vifm}_{i\textit{SD}r}$)
and public consumption ($\textit{vifm}_{iGr}$). The accounting
identity for these flows is thus:
$$\textit{vim}_{ir} = \sum_j \textit{vifm}_{ijr} +
\textit{vifm}_{iCr} + \textit{vifm}_{i\textit{SD}r} + \textit{vifm}_{iGr}\,.$$

Inputs to production of good $i$ ($Y_{ir}$) include intermediate
inputs (domestic $\textit{vdfm}_{ijr}$ and imported
$\textit{vifm}_{ijr}$), mobile factors of production
($\textit{vfm}_{m_{f},ir}$, where $m_{f}$ is a subset of the set $f$
designating all factors of production), and sector-specific factors of
production ($\textit{vfm}_{s_{f},ir}, s_f \subset f$). Factor market
equilibrium is given by an identity relating the value of factor
payments to factor income:
$$\sum_i \textit{vfm}_{fir} = \textit{evom}_{fr}\,,$$ and factor
earnings accrue to households.

International market clearance conditions require that region $r$
exports of good $i$ ($\textit{vxmd}_{ir}$ at the top of the figure)
equal the imports of the same good from the same region summed across
all trading partners ($\textit{vxmd}_{isr}$ at the bottom of the
figure): $$\textit{vxm}_{ir} = \sum_s \textit{vxmd}_{irs}\,,$$ where
$s$, an alias for $r$, indexes regions. Likewise, market clearance
conditions apply for international transportation services. The
supply-demand balance in the market for transportation service $j$
requires that the sum across all regions of service exports
($\textit{vst}_{ir}$, at the top of the figure) equals the sum across
all bilateral trade flows of service inputs ($\textit{vtwr}_{jisr}$ at
the bottom of the figure):
$$\sum_r \textit{vst}_{jr} = \sum_{isr} \textit{vtwr}_{jisr}$$

Turning to tax revenues and transfers, shown as dotted lines in figure
\ref{fig:region}, flows labeled with $\mathcal{R}$ correspond to tax
revenues. For each country, tax flows consist of indirect taxes on
production/exports of each good ($\mathcal{R}^Y_{ir}$), on consumption
($\mathcal{R}^C_r$), on public demand ($\mathcal{R}^G_r$) and on
imports ($\mathcal{R}^M_{ir}$). The regional budget constraint thus
relates tax payments ($\mathcal{R}^Y_{ir}, \mathcal{R}^C_r,
\mathcal{R}^G_r, \mathcal{R}^M_{ir}$), factor income
($\textit{evom}_{fr}$), and the current account deficit (i.e., net
transfers from abroad, $\textit{vb}_r$) to total private consumption
expenditure $\textit{vom}_{Cr}$, total public consumption expenditure
$\textit{vom}_{Gr}$, and total investment $\textit{vom}_{Ir}$,
yielding:
$$\textit{vom}_{Cr} + \textit{vom}_{Gr} + \textit{vom}_{Ir} = \sum_f
\textit{evom}_{fr} + \sum_i \mathcal{R}^Y_{ir} + \mathcal{R}^C_r +
\mathcal{R}^G_r + \sum_i \mathcal{R}^M_{ir} + \textit{vb}_r\,. $$

To this point we have outlined two types of consistency conditions
which are part of any \GTAP Data Base: market clearance (supply = demand
for all goods and factors), and income balance (net income = net
expenditure). A third set of identities involve net operating profits
by all sectors in the economy. In the core \GTAP model ``production''
takes place under conditions of perfect competition with constant
returns to scale, hence there are no excess profits, and the cost of
inputs must equal the value of outputs. This condition applies for
each production sector:

\begin{itemize}
\item[$Y_{ir}$:] $\sum_f \textit{vfm}_{fir} + \sum_j
\left(\textit{vdfm}_{ijr} + \textit{vifm}_{ijr} \right) +
\mathcal{R}^Y_{ir} = \textit{vom}_{ir}\,,$
\item[$C_{r}$:] $\sum_i \left(\textit{vdfm}_{iCr} +
\textit{vifm}_{iCr} \right) + \mathcal{R}^C_{ir} = \textit{vom}_{Cr}\,,$
\item[$I_{r}$:] $\sum_i \textit{vdfm}_{iIr} = \textit{vom}_{Ir}$
\item[$G_{r}$:] $\sum_i \left(\textit{vdfm}_{iGr} +
\textit{vifm}_{iGr} \right) + \mathcal{R}^G_{ir} =
\textit{vom}_{Gr}\,,$
\item[$M_{ir}$:] $\sum_s \left(\textit{vxmd}_{isr} + \sum_j
\textit{vtwr}_{jisr} \right) + \mathcal{R}^M_{ir} = \textit{vim}_{ir}$
\item[$FT_{fr}$:] $\textit{evom}_{fr} = \sum_i \textit{vfm}_{fir}$
\end{itemize}

\subsection[Optimization problems]{Decentralized optimization problems}

The benchmark identities presented in the previous section indicate
the market clearance, zero profit and income balance conditions which
define the \GTAP model. The displayed equations do not, however,
characterize the behavior of agents in the model. In a competitive
equilibrium setting, the standard assumption of optimizing atomistic
agents applies for both producers and consumers. This section lays out
the optimization problem of each component in the model, and thereby
provides the structure of production technology (production functions)
and preferences (characterizing final demand), as well as the
representation of trade. We also highlight where conceptual
differences between \GMR and \SOE models intervene.

Note that in order to simplify notation, we denote decision variables
corresponding to the benchmark data structures with the initial ``v''
replaced by ``d.'' Hence, while $vdfm_{jir}$ represents benchmark data
on intermediate demand for good $j$ in the production of good $i$ in
region $r$, $ddfm_{jir}$ represents the corresponding decision
variable in the equilibrium model. This approach to the scaling of
variables is consistent with the \GAMS code, and it provides a
flexible and transparent approach with respect to the calibration of
activity variables.

\subsubsection{Production technology}

Starting with producers, profit maximization in the constant returns
to scale setting is equivalent to cost minimization subject to
technical constraints. For sector $Y_{ir}$ we characterize input
choices as though they arose from minimization of unit costs:
\begin{eqnarray}\label{profitmax}
\min_{\textit{ddfm,dfm,difm}} \quad & c^D_{ir} + c^M_{ir} +
c^F_{ir} & \\
&s.t.&c^D_{ir} = \sum_j p^Y_{jr} (1 + t^{fd}_{jir})
\textit{ddfm}_{jir} \nonumber\\
&&c^M_{ir} = \sum_j p^M_{jr} (1 + t^{fi}_{jir}) \textit{difm}_{jir}\nonumber\\
&&c^F_{ir} = \sum_f (p^F_{fr}|_{f \in m_f} + p^S_{fir}|_{f \in s_f}) (1 +
t^f_{fir} ) \textit{dfm}_{fir}\nonumber\\
&&F_{ir}(\textit{ddfm,difm,dfm}) = Y_{ir}\nonumber
\end{eqnarray}
where $F(\cdot)$ represents the production function, which is
described by a nested constant-elasticity-of-substitution (CES) form,
with structure displayed in Figure \ref{prodfunY}.

\begin{figure}[t!]
\centering
\includegraphics{figures/prodfunY.tikz}
\vspace{0.1in}
\caption{CES nesting structure for production function $Y_{ir}=F_{ir}(\textit{ddfm,difm,dfm})$.}
\label{prodfunY}
\vspace{0.3in}
\end{figure}

In the figure, $\sigma$ values in different nests represent
substitution elasticities between inputs, with $\sigma^D_i =
\textit{esubd}_i$ measuring substitution possibility between
intermediate inputs produced domestically and imported from abroad
which are similarly a composite of imports from varieties from
different regions.  $\textit{esubva}_i$ represents the elasticity of
substitution between primary inputs in the value added nest. (Both of
these parameters are provided in the \GTAP 12 Data Base
\citep[see][]{HertelvanderMensbrugghe2015}.) Note further that the
specific source of tax revenue is indicated in this figure, consisting
of output taxes, taxes on intermediate inputs and taxes on factor
demands, all of which are applied on an ad-valorem basis.

One important difference between \GMR and \SOE models
occurs in the differentiation of output for domestic and export
markets. In particular, the supply of goods to domestic and export
markets are portrayed as arising from the following
profit-maximization problem:
\begin{eqnarray}\label{cet_dx}
\max_{\textit{dxm},\textit{ddm}} \quad && p^D_{ir} \textit{ddm}_{ir}
+ p^X_{ir} \textit{dxm}_{ir}
\\
&s.t.& G_{ir}(\textit{ddm}_{ir},\textit{dxm}_{ir}) = Y_{ir}
\end{eqnarray}
where $G$ is the \CET function with structure illustrated in Figure
\ref{prodfunDX}. In the \GMR model,
transformation elasticities $\eta^{DX}_{ir} = \textit{etrndx}_{i}$ are
set to infinity, which is the default value in the \GTAP 12 Data Base.
This implies that the supply price of output $p^Y_{jr}$ is the same if
output is used locally or in a different region. By contrast, in the
\SOE model we have that $\textit{etrndx}_{i} < \infty$,
which implies that output price may differ when supplied domestically
or abroad. In this case, $p^Y_{jr}$ in equation \eqref{profitmax} is
replaced by $p^D_{jr}.$ We come back to the choice of
$\textit{etrndx}_{i}$ in the \textsc{soe} model when we present our
illustrative calculations.

\begin{figure}[t!]
\centering
\includegraphics{figures/prodfunDX.tikz}
\vspace{0.1in}
\caption{CET transformation between domestic and export markets
$\textit{Y}_{ir} = G_{ir}(\textit{dxm},\textit{ddm})$.}
\label{prodfunDX}
\vspace{0.3in}
\end{figure}

\subsubsection{Preferences and final demand}

Private consumption consistent with utility maximization is portrayed
by minimization of the cost of a given level of aggregate consumption:
\begin{eqnarray}\label{utils}
\min_{\textit{ddfm}_{iCr},\textit{difm}_{iCr}} \quad &&\sum_i p^Y_{ir}
(1 + t^{pd}_{ir}) \textit{ddfm}_{iCr} + p^M_{ir} (1 +t^{pi}_{ir})
\textit{difm}_{iCr} \\
&s.t.& H_r(\textit{ddfm}_{iCr},\textit{difm}_{iCr}) = C_{ir} \nonumber
\end{eqnarray}
where $H_r$ represents final demand from the representative consumer.

Final demand in the core model is characterized by Cobb-Douglas
preferences. Alternative specifications included in the model logic
include both the \LES and \CDE expenditure.\footnote{The \LES model
involves subsistence and discretionary expenditure, which are
represented through composite price indices representing \CES
aggregates of both domestic and imported inputs, \texttt{P("DD",r)}
and \texttt{P("SD",r)}, respectively.  The \CDE model requires the
introduction of trade aggregation activities (\texttt{A(i,r)}) and
associated Armington composite price indices, \texttt{PA(i,r)}.} The
nested discretionary and subsistence demand functions are is displayed
in Figure \ref{prodfunC}.

\begin{figure}[t!]
\centering
\includegraphics{figures/prodfunC.tikz}
\vspace{0.1in}
\caption{LES nesting for discretionary and subsistence consumption.}
\label{prodfunC}
\vspace{0.3in}
\end{figure}

\subsubsection{Government and public consumption}

Public consumption in the model is represented as a fixed coefficient
(Leontief) aggregation of domestic-import composites. This formulation
introduces substitution at the second level between domestic and
imported inputs while holding sectoral commodity aggregates constant.
Figure \ref{prodfunG} illustrates the functional form.

\begin{figure}[t!]
\centering
\includegraphics{figures/prodfunG.tikz}
\vspace{0.1in}
\caption{Nesting structure for public consumption $G_r = G_r(\textit{ddfm}_{iGr},\textit{difm}_{iGr})$.}
\label{prodfunG}
\vspace{0.3in}
\end{figure}

\subsubsection{International trade}

The choice among imports from different trading partners is based on
Armington's idea of regionally differentiated products. The following
cost minimization problem formalizes this choice:
\begin{eqnarray}\label{optA}
\min_{\textit{dxmd,dtwr}} \quad &&\sum_s (1+t^{ms}_{isr}) \left(p^Y_{is}
(1-t^{xs}_{isr}) \textit{dxmd}_{isr} +
\sum_j p^T_j \textit{dtwr}_{jisr} \right)\\
&s.t.&A_{ir}(\textit{dxmd,dtwr}) = M_{ir} \nonumber
\end{eqnarray}
where $A$ is the import aggregation function, described by the nested
\CES-Leontief function shown in Figure \ref{prodfunA}. In
the case of the \SOE model, in which $\textit{etrndx}_{i}<\infty$,
$p^Y_{is}$ in (\ref{optA}) is replaced by $p^X_{is}$.

\begin{figure}[t!]
\centering
\includegraphics{figures/prodfunA.tikz}
\vspace{0.1in}
\caption{Armington aggregation of traded goods $M_{ir} = A_{ir}(\textit{dxmd,dtwr})$.}
\label{prodfunA}
\vspace{0.3in}
\end{figure}

Note that transportation services enter on a proportional basis with
imports from different countries, reflecting differences in unit
transportation margins across different goods and trading partners.
Therefore, substitution at the top level in an Armington composite
involves trading off of both imported goods and associated
transportation services. Trade flows are subject to export subsidies
and import tariffs, with subsidies paid by government in the exporting
region, and tariffs collected by government in the importing region.

The provision of international transportation services is modeled
through an aggregation of transportation services exported from
countries throughout the world. More specifically, we consider the
following cost minimization problem for the aggregation of
transportation services:
$$\min_{\textit{dst}} \quad \sum_r p^Y_{ir} \textit{dst}_{ir} \quad
\text{s.t.} \quad T_i(\textit{dst}) = \textit{YT}_i$$
where the aggregation function $T_i$ combines transport service
exports from multiple regions. The functional form which aggregates
services from different regions is illustrated in Figure
\ref{prodfunYT}.

\begin{figure}[t!]
\centering
\includegraphics{figures/prodfunYT.tikz}
\vspace{0.1in}
\caption{International transportation services aggregator
$\textit{YT}_j = T_j(\textit{dst})$.}
\label{prodfunYT}
\vspace{0.3in}
\end{figure}

\subsubsection{Supply of sector-specific factors}

Land and natural resources are portrayed as sector-specific factors of
production. These are supplied through
constant-elasticity-of-transformation (CET) production function which
allocates composite factors to sectoral markets. Formally, the supply
of sectoral factors of production is modeled through the following
profit-maximization problem:
\begin{eqnarray}\label{cet_input}
\max_{\textit{dfm}} \quad && \sum_j \textit{dfm}_{s_{f},jr} p^S_{s_{f},jr}\\
&s.t.& \Gamma_{s_{f},r}(\textit{dfm}) = \textit{evom}_{s_{f},r}
\end{eqnarray}
where $\Gamma$ is the \CET function with structure illustrated in
Figure \ref{prodfunFT}. Note that in the figure $\eta$ represents
transformation elasticities provided in the \GTAP 12 Data Base.

\begin{figure}[t!]
\centering
\includegraphics{figures/prodfunFT.tikz}
\vspace{0.1in}
\caption{Sector-specific factor \CET transformation function
$\textit{FT}_{s_{f}} = \Gamma_{s_{f}}(\textit{dfm})$.}
\label{prodfunFT}
\vspace{0.3in}
\end{figure}

%%% Local Variables:
%%% mode: LaTeX
%%% TeX-master: "documentation"
%%% End:
